\PassOptionsToPackage{dvipsnames,svgnames,table}{xcolor}
\documentclass[10pt]{article}
\usepackage[a4paper,margin=22mm]{geometry}
\usepackage[no-math]{fontspec}\setmainfont[SmallCapsFont={Latin Modern Roman Caps}]{Latin Modern Roman}
\usepackage{amsmath,amssymb,amsthm,mathtools,graphicx,xcolor,hyperref,xurl,xspace,ifthen,enumerate,textcomp}
\usepackage[shortlabels]{enumitem}
\setlength{\emergencystretch}{4em}
\newtheorem{theorem}{Theorem}[section]
\newtheorem{theo}[theorem]{Theorem}
\newtheorem{thm}[theorem]{Theorem}
\newtheorem{lemma}[theorem]{Lemma}
\newtheorem{lemm}[theorem]{Lemma}
\newtheorem{prop}[theorem]{Proposition}
\newtheorem{coro}[theorem]{Corollary}
\newtheorem{corollary}[theorem]{Corollary}
\newtheorem{fact}[theorem]{Fact}
\newtheorem{claim}[theorem]{Claim}
\theoremstyle{definition}
\newtheorem{defi}[theorem]{Definition}
\newtheorem{definition}[theorem]{Definition}
\newtheorem{rema}[theorem]{Remark}
\newtheorem{remark}[theorem]{Remark}
\newtheorem{exam}[theorem]{Example}
\newtheorem{example}[theorem]{Example}
\newenvironment{enonce*}[2][]{\par\medskip\noindent\textbf{#2.}\itshape}{\par\medskip}
\providecommand{\eg}{e.g.\ }
\providecommand{\ie}{i.e.\ }
\providecommand{\resp}{resp.\ }
\providecommand{\cf}{cf.\ }
\providecommand{\longthanks}[1]{\par\smallskip\noindent #1\par\smallskip}
\newtheorem{proposition}[theorem]{Proposition}
\newtheorem{conjecture}[theorem]{Conjecture}
\numberwithin{equation}{section}
\newcommand{\R}{{\mathbb R}}
\newcommand{\Z}{{\mathbb Z}}
\newcommand{\N}{{\mathbb N}}
\newcommand{\expref}[2]{\texorpdfstring{\hyperref[#2]{#1~\ref{#2}}}{#1~\ref{#2}}}
\newcommand{\epfmtdoi}[1]{\href{https://doi.org/#1}{doi:#1}}
\newcommand{\epfmt}[2]{#1:\,#2}
\newcommand{\AND}     {\text{\rm AND}}
\newcommand{\OR}      {\text{\rm OR}}
\newcommand{\PARITY}  {\text{\rm PARITY}}

%
 \DeclareMathOperator*{\Exp}{\mathbf{E}}
 \DeclareMathOperator*{\Prob}{\mathbf{P}}
 \newcommand{\adeg}{\deg_{1/3}}
 \DeclareMathOperator{\rdeg}{rdeg}
 \DeclareMathOperator{\rdegeps}{\rdeg_\epsilon}
 \newcommand{\degeps}{\deg_\epsilon}
 \newcommand{\degthr}{\deg_{\pm}}
 \DeclareMathOperator{\supp}{supp}


\newcommand{\1}{\mathbf{1}}
\newcommand{\tr}{^{\sf T}}
\newcommand{\eps}{\epsilon}
\newcommand{\e}{\mathrm{e}}  %
\newcommand{\moo}{\{-1,+1\}}
\newcommand{\moon}{\moo^n}
\newcommand{\mook}{\moo^k}
\newcommand{\moot}{\moo^t}
\newcommand{\moom}{\moo^m}
\newcommand{\zoo}{\{0,1\}}
\newcommand{\zoon}{\zoo^n}
\newcommand{\zook}{\zoo^k}
\newcommand{\zoot}{\zoo^t}
\newcommand{\zoom}{\zoo^m}
\newcommand{\zodn}{\{0,1,\dots,n\}}
\newcommand{\oneton}{\{1,2,\dots,n\}}

\newcommand{\tildeg}{\deg\hspace{-18pt}\widetilde{\rule{0mm}{6pt}\hspace{6mm}}}
\begin{document}
\begingroup\raggedright
\section*{1560. Tight approximate-degree lower bound for rectangular AND-OR trees}
Alexander A. Sherstov. \emph{Approximating the AND-OR Tree}. Theory of Computing volume 9 article 20 (2013), official complete journal PDF acquired 2026-09-30; canonical DOI 10.4086/toc.2013.v009a020 verified against official landing and printed PDF header. \url{https://theoryofcomputing.org/articles/v009a020/}. CC BY3.0.
\paragraph{Selected problem and scope (editorial).} Exact original Main Theorem1.1 only: the m-by-n rectangular AND-OR function has approximate degree Omega(sqrt(mn)). The selected outcome is the original lower bound; the cited tight upper bound is identified but not claimed as a new source-supplied proof. Complete approximate-degree/constant-error model, exact duality and Nisan-Szegedy input, source-supplied one-sided OR dual-witness proof, new signed-part averaging operator, full outer-AND approximation and block-degree contraction claims, and primary closing are retained. Linear-programming duality, ordinary polynomial multilinearity/linearity, probability conditioning and the named single-level AND/OR degree theorem are prerequisites. Witness/claims and the equivalent Theorem3.1 restatement remain uncounted core; communication corollaries and balanced parameter specializations are unselected.
\paragraph{Difficulty (editorial).} After the stated single-level degree and duality inputs, the proof still has to convert a one-sided dual witness into positive and negative conditional measures whose product averaging both approximates the outer AND and kills every low-degree block contribution. The two complete source-authored claims at439–488 establish these compatible approximation and degree-contraction requirements; applying only the known single-level AND/OR bounds would not give the rectangular composition bound. This is one H2 primal-dual composition construction above routine polynomial duality or interpolation qualification exercises.
\paragraph{Formatting (editorial).} Exact original human mathematical body and order retained; standalone typography, counter restoration and declared noncore printed locators only. Original comment prose excluded with percent guards retained; source typography and counters restored without upstream class execution. No mathematical repair or reconstruction.
\paragraph{Original source quirks (editorial).} The original approximate-degree notation uses the author-defined tildeg spacing macro in the introduction statement and adeg in the equivalent3.1 restatement; both exact source definitions and original statements are retained. Source constants, strict/nonstrict inequalities and signed-part conditioning are left unchanged. The full primary closing precedes the two complete claim proofs in the original author order; no reordering or mathematical repair is made.
\endgroup
\clearpage
\setcounter{section}{1}\setcounter{theorem}{0}
\begin{theorem}[Main result]
The function $f(x)=\bigwedge_{i=1}^m\bigvee_{j=1}^nx_{ij}$ has
approximate degree
\begin{align*}
\tildeg(f)=\Omega(\sqrt{mn})\,.
\end{align*}
%
\end{theorem}

This lower bound is tight for all $m$ and $n,$
by the results of H\o yer~et al.~\cite{hoyer-mosca-dewolf03and-or-tree}. 
\setcounter{section}{1}
\section{Preliminaries}
 
For a function $f\colon X\to\R$ on a finite set 
$X,$ we let $\|f\|_{\infty}=\max_{x\in X}|f(x)|$.
The total degree of a multivariate real polynomial $p\colon\R^n\to\R$
is denoted $\deg p$. We use the terms \emph{degree} and \emph{total
degree} interchangeably in this paper. For a function $f\colon X\to\R$
on a finite set $X\subset\R^n$, the \emph{$\epsilon$-approximate
degree} $\degeps(f)$ of $f$ is defined as the least degree of a
real polynomial $p$ with $\|f-p\|_{\infty}\leq\epsilon$. Throughout
this paper, we will work with the $\epsilon$-approximate degree for
a small constant $\epsilon>0$. For Boolean functions $f\colon X\to\zoo$,
the choice of constant $0<\epsilon<1/2$ affects the quantity $\degeps(f)$
by at most a constant factor:
\begin{align}
c\adeg(f)\leq\degeps(f)\leq C\adeg(f)\,,\label{eq:error-boosting}
\end{align}
where $c=c(\epsilon)$ and $C=C(\epsilon)$ are positive constants.
By convention, one studies $\epsilon=1/3$ as the canonical case and
reserves for it the special symbol $\tildeg(f)=\adeg(f)$. A dual
characterization~\cite{sherstov07quantum,sherstov11quantum-sdpt}
of the approximate degree is as follows.

\begin{fact}
\label{fact:dual}
Let $f\colon X\to\R$ be given, for a finite set
$X\subset\R^n$. Then $\degeps(f)\geq d$ if and only if there
exists a function $\psi\colon X\to\R$ such that
\begin{align*}
& \sum_{x\in X}|\psi(x)|=1\,,\\
& \sum_{x\in X}\psi(x)f(x)>\epsilon\,,
\intertext{and}
&\sum_{x\in X}\psi(x)p(x)=0
\end{align*}
for every polynomial $p$ of degree less than $d$.
\end{fact}

We adopt the usual definitions of the Boolean functions 
$\AND_n,\OR_n\colon\zoon\to\zoo$.
Their approximate degree was determined by Nisan 
and Szegedy~\cite{nisan-szegedy94degree}.

\begin{theorem}[Nisan and Szegedy]
\label{thm:nisan-szegedy}The functions $\AND_n$ and $\OR_n$
obey
\begin{align*}
\adeg(\AND_n)=\adeg(\OR_n)=\Theta(\sqrt{n})\,.
\end{align*}
\end{theorem}

By combining the above two theorems, Gavinsky and the author~\cite[Thm.~5.1]{GS09npconp}
obtained the following result, which plays a key role in this paper.

\begin{theorem}[Gavinsky and Sherstov]
\label{thm:dual}Fix any constant $0<\epsilon<1$. Then there exists
a constant $\delta=\delta(\epsilon)>0$ and a real function 
$\psi\colon\zoon\to\R$ such that
\begin{align}
& \sum_{x\in\zoon}|\psi(x)|=1\,,\label{eq:psi-bounded}\\
& \psi(0,0,\dots,0)<-\frac{1-\epsilon}2\,,\label{eq:psi-metric}
\intertext{and}
&\sum_{x\in\zoon}\psi(x)p(x)  =0\label{eq:psi-orthog}
\end{align}
for every polynomial $p$ of degree less than $\delta\sqrt{n}$.
\end{theorem}

For the sake of completeness, we include the proof.

\begin{proof}[Proof of \expref{Theorem}{thm:dual} 
(adapted from~\cite{GS09npconp})]
Recall from \expref{Theorem}{thm:nisan-szegedy} that $\adeg(\OR_n)=\Omega(\sqrt{n})$.
Therefore, \eqref{eq:error-boosting} shows that 
$\deg_{\frac{1-\epsilon}2}(\OR_n)\geq\delta\sqrt{n}$
for a sufficiently small constant $\delta=\delta(\epsilon)>0$. Now
the dual characterization of the approximate degree (\expref{Fact}{fact:dual})
provides a function $\psi\colon\zoon\to\R$ that obeys \eqref{eq:psi-bounded},
\eqref{eq:psi-orthog}, and
\begin{align}
\sum_{x\in\zoon}\psi(x)\OR_n(x)>\frac{1-\epsilon}2\,.\label{eq:correl-with-OR}
\end{align}
It remains to verify \eqref{eq:psi-metric}:
\begin{align*}
\psi(0,0,\dots,0) & =\sum_{x\in\zoon}\psi(x)(1-\OR_n(x))\\
& =-\sum_{x\in\zoon}\psi(x)\OR_n(x) & & \text{by \eqref{eq:psi-orthog}}\\
& <-\frac{1-\epsilon}2 & & \text{by
\eqref{eq:correl-with-OR}.}\qedhere
\end{align*}
\end{proof}

For probability distributions $\mu$ and $\lambda$ on finite sets
$X$ and $Y$, respectively, we let $\mu\times\lambda$ denote the
probability distribution on $X\times Y$ given by 
$(\mu\times\lambda)(x,y)=\mu(x)\lambda(y)$.
The \emph{support }of a probability distribution $\mu$ is defined
to be $\supp\mu=\{x:\mu(x)>0\}$.
 
 
\section{Main Result}
 
We are now in a position to prove our main result.

%
%
\begin{theorem}
The Boolean function $f(x)=\bigwedge_{i=1}^m\bigvee_{j=1}^nx_{ij}$
obeys
\begin{align}
\adeg(f)=\Omega(\sqrt{mn})\,.\label{eq:f-adeg-main-equation}
\end{align}
%
\end{theorem}

\begin{proof}
Let $\epsilon$ be an absolute constant to be named later, $0<\epsilon<1$.
Then by \expref{Theorem}{thm:dual}, there exists a constant $\delta=\delta(\epsilon)>0$
and a function $\psi\colon\zoon\to\R$ that
obeys~\eqref{eq:psi-bounded}--\eqref{eq:psi-orthog}.
Let $\mu$ be the probability distribution on $\zoon$ given by $\mu(x)=|\psi(x)|$.
Let $\mu_0$ and $\mu_1$ be the probability distributions induced
by $\mu$ on the sets $\{x:\psi(x)<0\}$ and $\{x:\psi(x)>0\},$ respectively.
Since $\sum_{x\in\zoon}\psi(x)=0,$ the sets $\{x:\psi(x)<0\}$ and
$\{x:\psi(x)>0\}$ are weighted equally by $\mu$. As a consequence,
\begin{align}
\mu & =\frac12\mu_1+\frac12\mu_0\,,\label{eq:mu-mu0-mu1}\\
\psi & =\frac12\mu_1-\frac12\mu_0\,.\label{eq:psi-mu0-mu1}
\end{align}
 
 
Consider the linear operator $L$ that maps functions $\phi\colon(\zoon)^m\to\R$
to functions $L\phi\colon\zoom\to\R$ according to
\begin{align*}
(L\phi)(z)=\Exp_{x_1\sim\mu_{z_1}}\cdots\Exp_{x_m\sim\mu_{z_m}}\phi(x_1,\dots,x_m)\,.
\end{align*}
Fix a real polynomial $p$ with
\begin{align}
\|f-p\|_{\infty}\leq\epsilon\,.\label{eq:F-p}
\end{align}
 
\begin{claim} \label{cla:Lf-AND}
\quad $\|\AND_m-Lf\|_{\infty}<\epsilon$.
\end{claim}
 
\begin{claim}
\label{cla:deg}
\quad $\deg p\geq\delta\sqrt{n}\;\deg Lp$.
\end{claim}

Before settling the claims, we finish the proof of the theorem. The
linearity of $L$ yields
\begin{align*}
\|\AND_m-Lp\|_{\infty} & \leq\underbrace{\|\AND_m-Lf\|_{\infty}}_{<\epsilon}
+\underbrace{\|L(f-p)\|_{\infty}}_{\leq\epsilon}<2\epsilon\,,
\end{align*}
where we have used~\eqref{eq:F-p} and \expref{Claim}{cla:Lf-AND} in
bounding the marked quantities. For $\epsilon=1/6,$ we arrive at
$\|\AND_m-Lp\|_{\infty}\leq1/3$ and therefore $\deg Lp=\Omega(\sqrt{m})$
by \expref{Theorem}{thm:nisan-szegedy}. Now \expref{Claim}{cla:deg} implies
that $\deg p=\Omega(\sqrt{mn})$.
\end{proof}
 
\begin{proof}[Proof of \expref{Claim}{cla:Lf-AND}]
By~\eqref{eq:psi-metric}, we have $\psi(x)>0$ only when $\OR_n(x)=1$.
Hence $\supp\mu_1\subseteq\OR_n^{-1}(1)$ and
\begin{align*}
(Lf)(1,1,\dots,1)
   =\Exp_{\mu_1\times\cdots\times\mu_1}[f]
   =\prod_{i=1}^m\Exp_{\mu_1}[\OR_n]=1\,.
\end{align*}
 
 
It remains to prove that $|(Lf)(z)|<\epsilon$ for every $z\ne(1,1,\dots,1)$.
We have
\begin{align*}
(Lf)(z)
  =\Exp_{\mu_{z_1}\times\cdots\times\mu_{z_m}}[f]
  =\prod_{i=1}^m\Exp_{\mu_{z_i}}[\OR_n]
  =\prod_{i=1}^m(1-\mu_{z_i}(0,0,\dots,0))\,,
\end{align*}
whence
\begin{align}
0\leq(Lf)(z)\leq1-\mu_0(0,0,\dots,0)\,.\label{eq:Lf-0-1-mu}
\end{align}
We know from~\eqref{eq:psi-metric} that $\psi(0,0,\dots,0)<-(1-\epsilon)/2,$
which means in particular that $(0,0,\dots,0)\in\supp\mu_0$. Therefore
\begin{align*}
\mu_0(0,0,\dots,0)
  =2\mu(0,0,\dots,0)
  =2|\psi(0,0,\dots,0)|
  >1-\epsilon\,,
\end{align*}
where the first step uses \eqref{eq:mu-mu0-mu1}. By~\eqref{eq:Lf-0-1-mu},
we conclude that $0\leq(Lf)(z)<\epsilon$.
\end{proof}
 
\begin{proof}[Proof of \expref{Claim}{cla:deg}]
By the linearity of $L$, it suffices to consider factored polynomials
$p$ of the form $p(x)=\prod_{i=1}^mp_i(x_{i,1},x_{i,2},\dots,x_{i,n})$.
In this case we have the convenient formula
\begin{align*}
(Lp)(z)=\prod_{i=1}^m\Exp_{\mu_{z_i}}[p_i]\,.
\end{align*}
By \eqref{eq:psi-orthog} and \eqref{eq:psi-mu0-mu1}, polynomials
$p_i$ of degree less than $\delta\sqrt{n}$ obey 
$\Exp_{\mu_0}[p_i]=\Exp_{\mu_1}[p_i]$
and therefore do not contribute to the degree of $Lp$. As a result,
\[
\deg Lp\leq|\{i:\deg p_i\geq\delta\sqrt{n}\}|
   \leq\frac{\deg p}{\delta\sqrt{n}}\,.\qedhere
\]
\end{proof}
\clearpage
\begingroup\raggedright
%
%
%
%
%
\providecommand{\bibhead}[1]{}
%
\expandafter\ifx\csname pdfbookmark\endcsname\relax%
  \providecommand{\tocrefpdfbookmark}{}
\else\providecommand{\tocrefpdfbookmark}{%
   \hypertarget{tocreferences}{}%
   \pdfbookmark[1]{References}{tocreferences}}%
\fi

\tocrefpdfbookmark
\begin{thebibliography}{10}

\bibitem{aaronson08tutorial}\bibhead{aaronson08tutorial}
{\sc Scott Aaronson}: The polynomial method in quantum and classical computing.
\newblock In {\em Proc. 49th FOCS}, p.~3. IEEE Comp. Soc. Press, 2008.
\newblock [\epfmtdoi{10.1109/FOCS.2008.91}]

\bibitem{aaronson-shi04distinctness}\bibhead{aaronson-shi04distinctness}
{\sc Scott Aaronson and Yaoyun Shi}: Quantum lower bounds for the collision and
  the element distinctness problems.
\newblock {\em J. ACM}, 51(4):595--605, 2004.
\newblock Preliminary version in
  \href{http://dx.doi.org/10.1145/509907.509999}{STOC'02}.
\newblock [\epfmtdoi{10.1145/1008731.1008735}]

\bibitem{ambainis05collision}\bibhead{ambainis05collision}
{\sc Andris Ambainis}: Polynomial degree and lower bounds in quantum
  complexity: {C}ollision and element distinctness with small range.
\newblock {\em Theory of Computing}, 1(3):37--46, 2005.
\newblock [\epfmtdoi{10.4086/toc.2005.v001a003}]

\bibitem{ACRSZ07nand}\bibhead{ACRSZ07nand}
{\sc Andris Ambainis, Andrew~M. Childs, Ben~W. Reichardt, Robert {\v S}palek,
  and Shengyu Zhang}: Any {AND-OR} formula of size {$N$} can be evaluated in
  time {$N^{1/2+o(1)}$} on a quantum computer.
\newblock {\em SIAM J. Comput.}, 39(6):2513--2530, 2010.
\newblock Preliminary version in
  \href{http://dx.doi.org/10.1109/FOCS.2007.57}{FOCS'07}.
\newblock [\epfmtdoi{10.1137/080712167}]

\bibitem{aspnes91voting}\bibhead{aspnes91voting}
{\sc James Aspnes, Richard Beigel, Merrick~L. Furst, and Steven Rudich}: The
  expressive power of voting polynomials.
\newblock {\em Combinatorica}, 14(2):135--148, 1994.
\newblock Preliminary version in
  \href{http://dx.doi.org/10.1145/103418.103461}{STOC'91}.
\newblock [\epfmtdoi{10.1007/BF01215346}]

\bibitem{beals-et-al01quantum-by-polynomials}\bibhead{beals-et-al01quantum-by-polynomials}
{\sc Robert Beals, Harry Buhrman, Richard Cleve, Michele Mosca, and Ronald~de
  Wolf}: Quantum lower bounds by polynomials.
\newblock {\em J. ACM}, 48(4):778--797, 2001.
\newblock Preliminary version in
  \href{http://dx.doi.org/10.1109/SFCS.1998.743485}{FOCS'98}.
\newblock [\epfmtdoi{10.1145/502090.502097}]

\bibitem{beigel93polynomial-method}\bibhead{beigel93polynomial-method}
{\sc Richard Beigel}: The polynomial method in circuit complexity.
\newblock In {\em Proc. 8th IEEE Conf. on Structure in Complexity Theory
  (SCT'93)}, pp. 82--95. IEEE Comp. Soc. Press, 1993.
\newblock [\epfmtdoi{10.1109/SCT.1993.336538}]

\bibitem{beigel94perceptrons}\bibhead{beigel94perceptrons}
{\sc Richard Beigel}: Perceptrons, {$\mathsf{PP}$}, and the polynomial
  hierarchy.
\newblock {\em Comput. Complexity}, 4(4):339--349, 1994.
\newblock Preliminary version in
  \href{http://dx.doi.org/10.1109/SCT.1992.215376}{SCT'92}.
\newblock [\epfmtdoi{10.1007/BF01263422}]

\bibitem{buhrman-et-al99small-error}\bibhead{buhrman-et-al99small-error}
{\sc Harry Buhrman, Richard Cleve, Ronald~de Wolf, and Christof Zalka}: Bounds
  for small-error and zero-error quantum algorithms.
\newblock In {\em Proc. 40th FOCS}, pp. 358--368. IEEE Comp. Soc. Press, 1999.
\newblock [\epfmtdoi{10.1109/SFFCS.1999.814607}]

\bibitem{buhrman-dewolf02DT-survey}\bibhead{buhrman-dewolf02DT-survey}
{\sc Harry Buhrman and Ronald de~Wolf}: Complexity measures and decision tree
  complexity: {A} survey.
\newblock {\em Theoret. Comput. Sci.}, 288(1):21--43, 2002.
\newblock [\epfmtdoi{10.1016/S0304-3975(01)00144-X}]

\bibitem{buhrman07pp-upp}\bibhead{buhrman07pp-upp}
{\sc Harry Buhrman, Nikolai~K. Vereshchagin, and Ronald de~Wolf}: On
  computation and communication with small bias.
\newblock In {\em Proc. 22nd IEEE Conf. on Computational Complexity (CCC'07)},
  pp. 24--32. IEEE Comp. Soc. Press, 2007.
\newblock [\epfmtdoi{10.1109/CCC.2007.18}]

\bibitem{buhrman-dewolf01polynomials}\bibhead{buhrman-dewolf01polynomials}
{\sc Harry Buhrman and Ronald~de Wolf}: Communication complexity lower bounds
  by polynomials.
\newblock In {\em Proc. 16th IEEE Conf. on Computational Complexity (CCC'01)},
  pp. 120--130. IEEE Comp. Soc. Press, 2001.
\newblock [\epfmtdoi{10.1109/CCC.2001.933879}]

\bibitem{bun-thaler13and-or}\bibhead{bun-thaler13and-or}
{\sc Mark Bun and Justin Thaler}: Dual lower bounds for approximate degree and
  {M}arkov-{B}ernstein inequalities.
\newblock In {\em Proc. 40th Internat. Colloq. on Automata, Languages and
  Programming (ICALP'13)}. Springer, 2013.
\newblock To appear. Preprint available at
  \href{http://arxiv.org/abs/1302.6191}{arXiv}.

\bibitem{GS09npconp}\bibhead{GS09npconp}
{\sc Dmitry Gavinsky and Alexander~A. Sherstov}: A separation of
  {$\mathsf{NP}$} and {$\mathsf{coNP}$} in multiparty communication complexity.
\newblock {\em Theory of Computing}, 6(1):227--245, 2010.
\newblock [\epfmtdoi{10.4086/toc.2010.v006a010}]

\bibitem{hoyer-mosca-dewolf03and-or-tree}\bibhead{hoyer-mosca-dewolf03and-or-tree}
{\sc Peter H{\o}yer, Michele Mosca, and Ronald~de Wolf}: Quantum search on
  bounded-error inputs.
\newblock In {\em Proc. 30th Internat. Colloq. on Automata, Languages and
  Programming (ICALP'03)}, pp. 291--299. Springer, 2003.
\newblock [\epfmtdoi{10.1007/3-540-45061-0\_25}]

\bibitem{kahn96incl-excl}\bibhead{kahn96incl-excl}
{\sc Jeff Kahn, Nathan Linial, and Alex Samorodnitsky}: Inclusion-exclusion:
  Exact and approximate.
\newblock {\em Combinatorica}, 16(4):465--477, 1996.
\newblock [\epfmtdoi{10.1007/BF01271266}]

\bibitem{KKMS}\bibhead{KKMS}
{\sc Adam~Tauman Kalai, Adam~R. Klivans, Yishay Mansour, and Rocco~A.
  Servedio}: Agnostically learning halfspaces.
\newblock {\em SIAM J. Comput.}, 37(6):1777--1805, 2008.
\newblock Preliminary version in
  \href{http://dx.doi.org/10.1109/SFCS.2005.13}{FOCS'05}.
\newblock [\epfmtdoi{10.1137/060649057}]

\bibitem{klauck07product-thms}\bibhead{klauck07product-thms}
{\sc Hartmut Klauck, Robert {\v S}palek, and Ronald~de Wolf}: Quantum and
  classical strong direct product theorems and optimal time-space tradeoffs.
\newblock {\em SIAM J. Comput.}, 36(5):1472--1493, 2007.
\newblock Preliminary version in
  \href{http://dx.doi.org/10.1109/FOCS.2004.52}{FOCS'04}.
\newblock [\epfmtdoi{10.1137/05063235X}]

\bibitem{KS01dnf}\bibhead{KS01dnf}
{\sc Adam~R. Klivans and Rocco~A. Servedio}: Learning {DNF} in time {$2^{\tilde
  O(n^{1/3})}$}.
\newblock {\em J. Comput. System Sci.}, 68(2):303--318, 2004.
\newblock Preliminary version in
  \href{http://dx.doi.org/10.1145/380752.380809}{STOC'01}.
\newblock [\epfmtdoi{10.1016/j.jcss.2003.07.007}]

\bibitem{krause94depth2mod}\bibhead{krause94depth2mod}
{\sc Matthias Krause and Pavel Pudl{\'a}k}: On the computational power of
  depth-2 circuits with threshold and modulo gates.
\newblock {\em Theoret. Comput. Sci.}, 174(1-2):137--156, 1997.
\newblock Preliminary version in
  \href{http://dx.doi.org/10.1145/195058.195103}{STOC'94}.
\newblock [\epfmtdoi{10.1016/S0304-3975(96)00019-9}]

\bibitem{KP98threshold}\bibhead{KP98threshold}
{\sc Matthias Krause and Pavel Pudl{\'a}k}: Computing {B}oolean functions by
  polynomials and threshold circuits.
\newblock {\em Comput. Complexity}, 7(4):346--370, 1998.
\newblock Preliminary version in
  \href{http://dx.doi.org/10.1109/SFCS.1995.492670}{FOCS'95}.
\newblock [\epfmtdoi{10.1007/s000370050015}]

\bibitem{linial-nisan90incl-excl}\bibhead{linial-nisan90incl-excl}
{\sc Nathan Linial and Noam Nisan}: Approximate inclusion-exclusion.
\newblock {\em Combinatorica}, 10(4):349--365, 1990.
\newblock Preliminary version in
  \href{http://dx.doi.org/10.1145/100216.100250}{STOC'90}.
\newblock [\epfmtdoi{10.1007/BF02128670}]

\bibitem{minsky88perceptrons}\bibhead{minsky88perceptrons}
{\sc Marvin~L. Minsky and Seymour~A. Papert}: {\em Perceptrons: {A}n
  Introduction to Computational Geometry}.
\newblock MIT Press, Cambridge, Mass., 1969.

\bibitem{muroga71threshold}\bibhead{muroga71threshold}
{\sc Saburo Muroga}: {\em Threshold Logic and Its Applications}.
\newblock John Wiley \& Sons, New York, 1971.

\bibitem{nisan-szegedy94degree}\bibhead{nisan-szegedy94degree}
{\sc Noam Nisan and Mario Szegedy}: On the degree of {B}oolean functions as
  real polynomials.
\newblock {\em Comput. Complexity}, 4(4):301--313, 1994.
\newblock Preliminary version in
  \href{http://dx.doi.org/10.1145/129712.129757}{STOC'92}.
\newblock [\epfmtdoi{10.1007/BF01263419}]

\bibitem{odonnell03degree}\bibhead{odonnell03degree}
{\sc Ryan O'Donnell and Rocco~A. Servedio}: New degree bounds for polynomial
  threshold functions.
\newblock {\em Combinatorica}, 30(3):327--358, 2010.
\newblock Preliminary version in
  \href{http://dx.doi.org/10.1145/780542.780592}{STOC'03}.
\newblock [\epfmtdoi{10.1007/s00493-010-2173-3}]

\bibitem{paturi-saks94rational}\bibhead{paturi-saks94rational}
{\sc Ramamohan Paturi and Michael~E. Saks}: Approximating threshold circuits by
  rational functions.
\newblock {\em Inform. and Comput.}, 112(2):257--272, 1994.
\newblock Preliminary version in
  \href{http://dx.doi.org/10.1109/FSCS.1990.89559}{FOCS'90}.
\newblock [\epfmtdoi{10.1006/inco.1994.1059}]

\bibitem{razborov02quantum}\bibhead{razborov02quantum}
{\sc Alexander~A. Razborov}: Quantum communication complexity of symmetric
  predicates.
\newblock {\em Izvestiya: Mathematics}, 67(1):145--159, 2003.
\newblock [\epfmtdoi{10.1070/IM2003v067n01ABEH000422}]

\bibitem{RS07dc-dnf}\bibhead{RS07dc-dnf}
{\sc Alexander~A. Razborov and Alexander~A. Sherstov}: The sign-rank of
  {$\mathsf{AC}^0$}.
\newblock {\em SIAM J. Comput.}, 39(5):1833--1855, 2010.
\newblock Preliminary version in
  \href{http://dx.doi.org/10.1109/FOCS.2008.19}{FOCS'08}.
\newblock [\epfmtdoi{10.1137/080744037}]

\bibitem{reichardt-reflections11soda}\bibhead{reichardt-reflections11soda}
{\sc Ben~W. Reichardt}: Reflections for quantum query algorithms.
\newblock In {\em Proc. 22nd Ann. ACM-SIAM Symp. on Discrete Algorithms
  (SODA'11)}, pp. 560--569. ACM Press, 2011.
\newblock
  \href{https://www.siam.org/proceedings/soda/2011/SODA11_044_reichardtb.pdf}{SIAM}.
\newblock [\epfmt{acm}{2133080}]

\bibitem{saks93slicing}\bibhead{saks93slicing}
{\sc Michael~E. Saks}: Slicing the hypercube.
\newblock In {\sc Keith Walker}, editor, {\em Surveys in Combinatorics, 1993},
  volume 187 of {\em London Mathematical Society Lecture Note Series}, pp.
  211--256. Cambridge Univ. Press, 1993.
\newblock [\epfmtdoi{10.1017/CBO9780511662089.009}]

\bibitem{dual-survey}\bibhead{dual-survey}
{\sc Alexander~A. Sherstov}: Communication lower bounds using dual polynomials.
\newblock {\em Bulletin of the {EATCS}}, 95:59--93, 2008.
\newblock \href{http://www.eatcs.org/images/bulletin/beatcs95.pdf}{EATCS}.

\bibitem{sherstov07inclexcl-ccc}\bibhead{sherstov07inclexcl-ccc}
{\sc Alexander~A. Sherstov}: Approximate inclusion-exclusion for arbitrary
  symmetric functions.
\newblock {\em Comput. Complexity}, 18(2):219--247, 2009.
\newblock Preliminary version in
  \href{http://dx.doi.org/10.1109/CCC.2008.18}{CCC'08}.
\newblock [\epfmtdoi{10.1007/s00037-009-0274-4}]

\bibitem{sherstov09hshs}\bibhead{sherstov09hshs}
{\sc Alexander~A. Sherstov}: The intersection of two halfspaces has high
  threshold degree.
\newblock In {\em Proc. 50th FOCS}, pp. 343--362. IEEE Comp. Soc. Press, 2009.
\newblock [\epfmtdoi{10.1109/FOCS.2009.18}]

\bibitem{sherstov07ac-majmaj}\bibhead{sherstov07ac-majmaj}
{\sc Alexander~A. Sherstov}: Separating {$\mathsf{AC}^0$} from depth-2 majority
  circuits.
\newblock {\em SIAM J. Comput.}, 38(6):2113--2129, 2009.
\newblock Preliminary version in
  \href{http://dx.doi.org/10.1145/1250790.1250834}{STOC'07}.
\newblock [\epfmtdoi{10.1137/08071421X}]

\bibitem{sherstov07quantum}\bibhead{sherstov07quantum}
{\sc Alexander~A. Sherstov}: The pattern matrix method.
\newblock {\em SIAM J. Comput.}, 40(6):1969--2000, 2011.
\newblock Preliminary version in
  \href{http://dx.doi.org/10.1145/1374376.1374392}{FOCS'08}.
\newblock [\epfmtdoi{10.1137/080733644}]

\bibitem{sherstov12mdisj}\bibhead{sherstov12mdisj}
{\sc Alexander~A. Sherstov}: The multiparty communication complexity of set
  disjointness.
\newblock In {\em Proc. 44th STOC}, pp. 525--548. ACM Press, 2012.
\newblock [\epfmtdoi{10.1145/2213977.2214026}]

\bibitem{sherstov11quantum-sdpt}\bibhead{sherstov11quantum-sdpt}
{\sc Alexander~A. Sherstov}: Strong direct product theorems for quantum
  communication and query complexity.
\newblock {\em SIAM J. Comput.}, 41(5):1122--1165, 2012.
\newblock Preliminary version in
  \href{http://dx.doi.org/10.1145/1993636.1993643}{STOC'11}.
\newblock [\epfmtdoi{10.1137/110842661}]

\bibitem{sherstov12noisy}\bibhead{sherstov12noisy}
{\sc Alexander~A. Sherstov}: Making polynomials robust to noise.
\newblock {\em Theory of Computing}, 9(18):593--615, 2013.
\newblock Preliminary version in
  \href{http://dx.doi.org/10.1145/2213977.2214044}{STOC'12}.
\newblock [\epfmtdoi{10.4086/toc.2013.v009a018}]

\bibitem{shi-linear}\bibhead{shi-linear}
{\sc Yaoyun Shi}: Approximating linear restrictions of {B}oolean functions.
\newblock Manuscript. Available at
  \href{http://web.eecs.umich.edu/~shiyy/mypapers/linear02-j.ps}{author's home
  page}, 2002.

\bibitem{siu-roy-kailath94rational}\bibhead{siu-roy-kailath94rational}
{\sc Kai-Yeung Siu, Vwani~P. Roychowdhury, and Thomas Kailath}: Rational
  approximation techniques for analysis of neural networks.
\newblock {\em IEEE Trans. Inform. Theory}, 40(2):455--466, 1994.
\newblock Preliminary version in
  \href{http://books.nips.cc/nips05.html}{NIPS'92}.
\newblock [\epfmtdoi{10.1109/18.312168}]

\bibitem{tt99DNF-incl-excl}\bibhead{tt99DNF-incl-excl}
{\sc Jun Tarui and Tatsuie Tsukiji}: Learning {DNF} by approximating
  inclusion-exclusion formulae.
\newblock In {\em Proc. 14th IEEE Conf. on Computational Complexity (CCC'99)},
  pp. 215--221. IEEE Comp. Soc. Press, 1999.
\newblock [\epfmtdoi{10.1109/CCC.1999.766279}]

\bibitem{de-wolf08approx-degree}\bibhead{de-wolf08approx-degree}
{\sc {Ronald de} Wolf}: A note on quantum algorithms and the minimal degree of
  $\epsilon$-error polynomials for symmetric functions.
\newblock {\em Quantum Inf. Comput.}, 8(10):943--950, 2008.
\newblock
  \href{http://www.rintonpress.com/journals/qiconline.html\#v8n10}{QIC}.
\newblock [\epfmt{acm}{2016989}]

\end{thebibliography}
\endgroup
\end{document}
